% ============================================================
%  CHAPTER 4 — EZZAOUIA WEB PLATFORM
%  Design, Implementation and Security
%  EZZAOUIA Platform — MARETAP S.A. — PFE 2025–2026
%  Authentication: Django Sessions + CSRF Tokens (NOT JWT)
% ============================================================

\chapter{EZZAOUIA Web Platform --- Design, Implementation and Security}
\label{chap:platform}

\setcounter{section}{0}
\setcounter{figure}{0}
\setcounter{table}{0}

% ─────────────────────────────────────────────────────────────
\section{Introduction}
\begin{sloppypar}

With the Data Warehouse fully designed, populated, and validated as described in
the previous chapter, the development focus shifted to the construction of the
EZZAOUIA web platform --- the central deliverable of this internship project.
The platform is a full-stack web application built on Django and React, designed
to serve as the unified interface through which MARETAP~S.A.'s engineers and
managers interact with the company's production data, analytical dashboards, and
AI-assisted tooling. Rather than replacing any existing operational system, the
EZZAOUIA platform layers on top of the Data Warehouse and exposes its content
through a secure, role-aware, and user-friendly interface that consolidates what
was previously scattered across Excel files, printed TCM reports, and manual
calculations.

This chapter covers the complete lifecycle of the platform's construction across
two development sprints. Sprint~2 focused on the Django backend: defining the
project's application architecture, modelling the database entities, exposing KPI
data through a REST API, implementing session-based authentication with CSRF
protection, and building the role-based access control system. Sprint~3 focused
on the React frontend: assembling the component hierarchy, implementing the
role-gated navigation sidebar, building the KPI dashboard with Chart.js,
embedding Power~BI reports via secure iframes, creating the user management
interface, and integrating the AI chatbot and file ingestion modules. The chapter
also details the platform's security architecture, its containerised deployment
via Docker Compose, and the key technical decisions that shaped its final design.

\end{sloppypar}

% ─────────────────────────────────────────────────────────────
\section{Sprint 2 User Stories --- Django Backend \& REST API}
\begin{sloppypar}

\textbf{Period:} March 2 -- March 29, 2026

\noindent The following user stories were identified and implemented during
Sprint~2 to cover the backend infrastructure, data access layer, authentication
system, and REST API of the EZZAOUIA platform.

\begin{table}[H]
\caption{Sprint 2 --- Django Backend \& REST API User Stories}
\centering
\renewcommand{\arraystretch}{1.5}
\begin{tabular}{|p{2cm}|p{2cm}|p{9.5cm}|}
\hline
\rowcolor{blue!15}
\textbf{ID} & \textbf{Role} & \textbf{User Story} \\ \hline
US-13 & Admin &
  As an admin, I want a Django project structured into multiple applications
  (\texttt{accounts}, \texttt{dashboard}, \texttt{chatbot}, \texttt{reports})
  so that the codebase is modular and maintainable. \\ \hline
\rowcolor{gray!10}
US-14 & Admin &
  As an admin, I want Django models (\texttt{Well}, \texttt{FactProduction},
  \texttt{FactTankLevel}) mapped to the SQL Server Data Warehouse via the Django
  ORM so that production data is queryable without writing raw SQL. \\ \hline
US-15 & User &
  As a user, I want a secure session-based authentication system (login and
  logout with CSRF protection) so that only authorised personnel can access the
  platform. \\ \hline
\rowcolor{gray!10}
US-16 & Admin &
  As an admin, I want a role-based access control system distinguishing two
  roles --- \textit{admin} and \textit{user} --- so that each profile sees only
  the features relevant to their responsibilities. \\ \hline
US-17 & User &
  As a user, I want REST API endpoints exposing KPIs (BOPD, BSW, GOR, top
  wells, monthly trend, tank levels) so that the React frontend can consume
  live production data. \\ \hline
\rowcolor{gray!10}
US-18 & Admin &
  As an admin, I want a user management API (create, list, update, and delete
  users; assign roles; generate temporary passwords) so that platform access can
  be administered without direct database interaction. \\ \hline
US-19 & Admin &
  As an admin, I want Power~BI report embed URLs stored in a Django model
  (\texttt{PowerBIReport}) with an auto-converting \texttt{embed\_src} property
  so that reports can be managed entirely from the admin interface. \\ \hline
\end{tabular}
\end{table}

% ── USE CASE DIAGRAMS — Sprint 2 ─────────────────────────────

\begin{figure}[H]
  \centering
  \includegraphics[width=14cm]{images/authentification.drawio.png}
  \caption{UC01 --- Authentication \& Session Management Use Case Diagram}
  \label{fig:uc01}
\end{figure}

\begin{figure}[H]
  \centering
  \includegraphics[width=14cm]{images/UC02_KPI_Dashboard_Overview.drawio.png}
  \caption{UC02 --- KPI Dashboard Overview Use Case Diagram}
  \label{fig:uc02}
\end{figure}

\begin{figure}[H]
  \centering
  \includegraphics[width=14cm]{images/User Management.drawio.png}
  \caption{UC03 --- User Management Use Case Diagram}
  \label{fig:uc03}
\end{figure}

\begin{figure}[H]
  \centering
  \includegraphics[width=14cm]{images/UC04_PowerBI_Report_Embedding.drawio.png}
  \caption{UC04 --- Power~BI Report Embedding Use Case Diagram}
  \label{fig:uc04}
\end{figure}

\end{sloppypar}

% ─────────────────────────────────────────────────────────────
\section{Sprint 3 User Stories --- React Frontend \& Platform Integration}
\begin{sloppypar}

\textbf{Period:} March 30 -- April 26, 2026

\begin{table}[H]
\caption{Sprint 3 --- React Frontend \& Platform Integration User Stories}
\centering
\renewcommand{\arraystretch}{1.5}
\begin{tabular}{|p{2cm}|p{2cm}|p{9.5cm}|}
\hline
\rowcolor{blue!15}
\textbf{ID} & \textbf{Role} & \textbf{User Story} \\ \hline
US-20 & User &
  As a user, I want a React single-page application with React Router so that
  navigation between modules is seamless without full page reloads. \\ \hline
\rowcolor{gray!10}
US-21 & User &
  As a user, I want a role-gated sidebar with pinned Profile and Logout buttons
  at the bottom so that navigation is intuitive and access-controlled by
  role. \\ \hline
US-22 & User &
  As a user, I want a KPI dashboard page displaying BOPD, BSW, GOR, top wells,
  and monthly production trends via Chart.js so that I can monitor field
  performance in real time. \\ \hline
\rowcolor{gray!10}
US-23 & User &
  As a user, I want Power~BI dashboards embedded directly inside the platform
  via secure iframes so that I can access all four reports without leaving the
  application. \\ \hline
US-24 & Admin &
  As an admin, I want a user management interface with a temporary password modal
  that persists until manually closed, and email failure warnings when SMTP
  delivery fails, so that user onboarding is reliable. \\ \hline
\rowcolor{gray!10}
US-25 & User &
  As a user, I want an AI-assisted chatbot powered by a RAG architecture
  (Ollama~+~ChromaDB) so that I can query production data and documents using
  natural language. \\ \hline
US-26 & Admin &
  As an admin, I want a file ingestion module backed by Celery asynchronous
  tasks so that uploaded documents are processed and indexed into ChromaDB
  without blocking the main application thread. \\ \hline
\rowcolor{gray!10}
US-27 & Admin &
  As an admin, I want the core platform services containerised with Docker
  Compose (\texttt{web}, \texttt{celery}, \texttt{celery-beat}, and
  \texttt{redis}) so that deployment is reproducible across
  environments. \\ \hline
US-28 & User &
  As a user, I want to view and edit my profile information and change my
  password so that my personal information stays up to date. \\ \hline
\end{tabular}
\end{table}

\begin{figure}[H]
  \centering
  \includegraphics[width=0.9\textwidth]{images/UC05_Chatbot_RAG.drawio.png}
  \caption{UC05 --- RAG Chatbot \& AI-Powered Analysis Use Case Diagram}
  \label{fig:uc05}
\end{figure}

\begin{figure}[H]
  \centering
  \includegraphics[width=0.9\textwidth]{images/UC06_File_Ingestion.drawio.png}
  \caption{UC06 --- File Ingestion Module Use Case Diagram}
  \label{fig:uc06}
\end{figure}

\begin{figure}[H]
  \centering
  \includegraphics[width=0.9\textwidth]{images/UC08_Profile_Management.drawio.png}
  \caption{UC08 --- Profile Management Use Case Diagram}
  \label{fig:uc08}
\end{figure}

\noindent\textit{Note: UC07 --- Production Forecasting was designed during
Sprint~3 but implemented in Sprint~4. Its use case diagram and full
implementation are presented in Chapter~6.}

\end{sloppypar}

% ─────────────────────────────────────────────────────────────
\section{Requirements Analysis}
\begin{sloppypar}

\subsection{Functional Requirements}

\begin{itemize}
  \item A secure, session-based authentication system enabling MARETAP personnel
    to log in with their credentials and receive role-specific access to the
    platform's modules, with CSRF token protection applied to every
    state-changing request.

  \item A role-based access control system distinguishing two profiles ---
    \textbf{admin} and \textbf{user} --- where administrators have access to
    all platform modules, while standard users are restricted to dashboards and
    reporting features.

  \item A KPI dashboard that dynamically queries the Data Warehouse through a
    REST API and renders production indicators using interactive Chart.js
    visualisations.

  \item An embedded Power~BI reporting module that renders four production
    dashboards directly within the platform interface via authenticated
    iframes.

  \item A user management interface allowing administrators to create accounts,
    assign roles, generate temporary passwords, and receive warnings when email
    delivery fails.

  \item An AI-assisted chatbot powered by a RAG architecture, enabling engineers
    to query production data and technical documents using natural language.

  \item A document ingestion module allowing administrators to upload files that
    are asynchronously processed and indexed into ChromaDB via Celery.

  \item A navigation sidebar that adapts its visible links based on the
    authenticated user's role.
\end{itemize}

\subsection{Non-Functional Requirements}

\begin{itemize}
  \item \textbf{Security:} All API endpoints are protected by Django's session
    authentication and CSRF token verification. The \texttt{sessionid} cookie is
    \texttt{HttpOnly}. CORS is strictly configured to allow only the React
    frontend's origin.

  \item \textbf{Performance:} The React SPA eliminates full page reloads.
    Celery offloads all file processing tasks from the main Django request
    thread.

  \item \textbf{Reliability:} Authentication errors, API failures, and SMTP
    delivery failures are handled gracefully through user-facing notifications.

  \item \textbf{Maintainability:} The Django project is split into
    \texttt{accounts}, \texttt{dashboard}, \texttt{chatbot}, and
    \texttt{reports} applications, each with a clearly defined
    responsibility.

  \item \textbf{Scalability:} Docker Compose allows each containerised service
    to be scaled independently.

  \item \textbf{Compatibility:} The platform runs entirely on MARETAP's
    internal network. The SQL Server Data Warehouse is accessed via ODBC
    Driver~18.

  \item \textbf{Usability:} All critical actions are reachable within two
    clicks from the authenticated landing page.
\end{itemize}

\end{sloppypar}

% ─────────────────────────────────────────────────────────────
\section{Technical Stack}
\begin{sloppypar}

\begin{table}[H]
\centering
\caption{EZZAOUIA Platform --- Full Technical Stack}
\label{tab:tech_stack}
\renewcommand{\arraystretch}{1.5}
\begin{tabular}{|l|l|p{7cm}|}
\hline
\rowcolor{blue!20}
\textbf{Layer} & \textbf{Technology} & \textbf{Role in the Platform} \\
\hline
Backend Framework & Django 4.x &
  REST API, ORM, authentication, admin interface, and business logic. \\
\hline
\rowcolor{gray!10}
API Layer & Django REST Framework &
  Serialisers, ViewSets, and session-protected endpoints. \\
\hline
Authentication & Django Sessions + CSRF &
  \texttt{sessionid} cookie (\texttt{HttpOnly}) and \texttt{csrftoken} read
  by Axios on each mutating request. \\
\hline
\rowcolor{gray!10}
Database (DWH) & Microsoft SQL Server &
  On-premise Data Warehouse storing all production facts and dimensions. \\
\hline
ORM Driver & pyodbc + ODBC Driver~18 &
  Low-level SQL Server connectivity from the Django ORM layer. \\
\hline
\rowcolor{gray!10}
Async Task Queue & Celery + Redis &
  Background document processing and file ingestion pipelines. \\
\hline
Frontend Framework & React~18 &
  Component-based single-page application. \\
\hline
\rowcolor{gray!10}
Routing & React Router~v6 &
  Client-side navigation with role-protected route guards. \\
\hline
Data Visualisation & Chart.js &
  Interactive KPI charts: bar, line, doughnut, and radar. \\
\hline
\rowcolor{gray!10}
BI Reporting & Power~BI (iframe embed) &
  Embedded analytical dashboards with \texttt{autoAuth=true}. \\
\hline
AI / LLM & Ollama &
  Local LLM inference for the RAG chatbot, running on the host machine. \\
\hline
\rowcolor{gray!10}
Vector Store & ChromaDB &
  Semantic search and document chunk retrieval for the RAG pipeline,
  persisted on the host filesystem. \\
\hline
Containerisation & Docker + Docker Compose &
  Four-service orchestration (\texttt{web}, \texttt{celery},
  \texttt{celery-beat}, \texttt{redis}) for reproducible deployment. \\
\hline
\rowcolor{gray!10}
Email & Microsoft~365 SMTP &
  Transactional email for temporary password delivery. \\
\hline
\end{tabular}
\end{table}

\end{sloppypar}

% ─────────────────────────────────────────────────────────────
\section{Platform Architecture}
\begin{sloppypar}

\subsection{General Architecture Overview}

The EZZAOUIA platform follows a decoupled client-server architecture. All
requests from the React frontend carry the Django session cookie
(\texttt{sessionid}) for authentication and a CSRF token in the
\texttt{X-CSRFToken} header for write-operation protection.

The platform is composed of \textbf{four} Docker containers orchestrated by
Docker Compose. Two additional services --- ChromaDB and Ollama --- run directly
on the host machine outside the Docker network and are accessible to the
containers via the host network interface:

\begin{itemize}
  \item \textbf{\texttt{web} container} --- port~8000. Hosts the Django REST
    API; handles session authentication, business logic, KPI computation, and
    user management.

  \item \textbf{\texttt{celery} container} --- no exposed port. Runs the Celery
    worker process; consumes document ingestion tasks dispatched by Django to
    the Redis message broker queue.

  \item \textbf{\texttt{celery-beat} container} --- no exposed port. Runs the
    Celery Beat periodic task scheduler, responsible for time-driven background
    jobs.

  \item \textbf{\texttt{redis} container} --- port~6379. Uses the official
    \texttt{redis:7-alpine} image; serves as the Celery message broker within
    the Docker Compose network.

  \item \textbf{ChromaDB (host process)} --- not containerised. Accessed as an
    embedded Python library within the \texttt{web} and \texttt{celery}
    containers; persists vector data to the host filesystem at
    \texttt{./chroma\_db/}.

  \item \textbf{Ollama (host process)} --- not containerised. Runs as a local
    inference server on the host machine (default:
    \texttt{http://127.0.0.1:11434}); queried by Django over the host
    network.
\end{itemize}

\begin{figure}[H]
    \centering
    \includegraphics[width=0.95\textwidth]{images/platform_architecture.png}
    \caption{General Architecture of the EZZAOUIA Platform}
    \label{fig:platform_arch}
\end{figure}

\subsection{Django Project Structure}

\begin{table}[H]
\centering
\caption{Django Application Structure of the EZZAOUIA Backend}
\label{tab:django_apps}
\renewcommand{\arraystretch}{1.5}
\begin{tabular}{|l|p{10cm}|}
\hline
\rowcolor{blue!20}
\textbf{Application} & \textbf{Responsibilities} \\
\hline
\texttt{accounts} &
  Custom user model, session-based login/logout views, role management, user
  CRUD API, temporary password generation, and SMTP email delivery. \\
\hline
\rowcolor{gray!10}
\texttt{dashboard} &
  Django ORM models mapped to the DWH, KPI computation logic, and REST API
  endpoints for all dashboard indicators. \\
\hline
\texttt{chatbot} &
  RAG pipeline orchestration: document ingestion, ChromaDB indexing, Ollama
  query handling, and conversational response assembly. \\
\hline
\rowcolor{gray!10}
\texttt{reports} &
  \texttt{PowerBIReport} model, embed URL auto-conversion, and API endpoint
  serving report metadata to the frontend. \\
\hline
\end{tabular}
\end{table}

\end{sloppypar}

% ─────────────────────────────────────────────────────────────
\section{Backend Implementation}
\begin{sloppypar}

\subsection{Data Models}

All Django models that reference DWH tables are configured with
\texttt{managed~=~False}, instructing Django not to create, alter, or drop these
tables during migrations. This ensures that the existing Data Warehouse schema
remains untouched by the platform's migration lifecycle.

% ── CLASS DIAGRAM ─────────────────────────────────────────────
\begin{figure}[H]
    \centering
    \fbox{\parbox{0.95\textwidth}{\centering\vspace{2cm}
      \textbf{[ CLASS DIAGRAM --- DJANGO ORM MODELS ]}\\[0.3cm]
      \texttt{CustomUser} $\cdot$ \texttt{Well} (\texttt{DimWell})
        $\cdot$ \texttt{FactProduction}\\
      \texttt{FactTankLevel} $\cdot$ \texttt{PowerBIReport}\\
      with attributes, relationships, and \texttt{managed=False} annotations\\[0.3cm]
      \textit{[ Insert class diagram image here ]}
    \vspace{2cm}}}
    \caption{Django ORM Model Relationships in the EZZAOUIA Backend}
    \label{fig:django_models}
\end{figure}

\subsubsection{Well Model}

Maps to \texttt{DimWell}. Exposes the full technical reference of each Ezzaouia
well, including production method, power type, and operational ordering.

\subsubsection{FactProduction Model}

Maps to the central \texttt{FactProduction} fact table. Carries
\texttt{DailyOilPerWellSTBD}, \texttt{DailyGasPerWellMSCF}, and
\texttt{WaterCutPercentage}, linked via foreign keys to \texttt{DimDate},
\texttt{DimWell}, and \texttt{DimWellStatus}. All KPI computations use
Django ORM's \texttt{annotate()}, \texttt{aggregate()}, and
\texttt{values()} methods.

\subsubsection{FactTankLevel Model}

Maps to \texttt{FactTankLevel}. Stores daily volume measurements per tank,
linked to \texttt{DimDate} and \texttt{DimTank}.

\subsubsection{PowerBIReport Model}

Stored in a Django-managed table (not part of the DWH schema). The
\texttt{embed\_src} Python property automatically converts a raw Power~BI URL
into the \texttt{reportEmbed} format with
\texttt{autoAuth=true\&ctid=\{tenant\_id\}}.

% ─────────────────────────────────────────────────────────────
\subsection{Authentication System}

\subsubsection{Session-Based Authentication}

The EZZAOUIA platform uses \textbf{Django's built-in session framework} as its
authentication mechanism --- the native, recommended approach for on-premise web
applications. This was chosen deliberately over token-based alternatives because
the platform is deployed entirely within MARETAP's internal network, where
server-side session management is both simpler to maintain and equally secure.

The authentication flow proceeds as follows:

\begin{enumerate}
  \item The user submits credentials via the React login form.
  \item Django validates the credentials against the \texttt{CustomUser} model
    and, upon success, creates a server-side session stored in the database.
  \item Django sets a \texttt{sessionid} cookie in the browser response. This
    cookie is \texttt{HttpOnly} --- inaccessible to JavaScript --- which
    protects against XSS-based session hijacking.
  \item On every subsequent API request, the browser automatically sends the
    \texttt{sessionid} cookie, which Django uses to identify the authenticated
    user.
  \item On logout, Django destroys the server-side session and instructs the
    browser to clear the cookie.
\end{enumerate}

The following session parameters are configured in \texttt{settings.py}:

\begin{itemize}
  \item \texttt{SESSION\_COOKIE\_AGE} --- controls the session lifetime in
    seconds.
  \item \texttt{SESSION\_EXPIRE\_AT\_BROWSER\_CLOSE} --- when set to
    \texttt{True}, the session expires when the browser window is closed.
  \item \texttt{SESSION\_COOKIE\_HTTPONLY = True} --- prevents JavaScript
    from accessing the session cookie.
  \item \texttt{SESSION\_COOKIE\_SECURE = True} --- ensures the cookie is
    transmitted only over HTTPS in production.
\end{itemize}

% ── SEQUENCE DIAGRAM: Session Auth ───────────────────────────
\begin{figure}[H]
  \centering
  \fbox{\parbox{0.9\textwidth}{\centering\vspace{1.5cm}
    \textbf{[ SEQUENCE DIAGRAM ]}\\[0.3cm]
    SD01 --- Django Session Authentication Flow\\[0.2cm]
    User $\rightarrow$ React Login Form
      $\rightarrow$ \texttt{POST /api/auth/login/} (+ \texttt{csrftoken})\\
    $\rightarrow$ Django validates credentials
      $\rightarrow$ Creates session (database)\\
    $\rightarrow$ Sets \texttt{sessionid} cookie (\texttt{HttpOnly})
      $\rightarrow$ Redirect to Dashboard\\[0.2cm]
    On logout: \texttt{POST /api/auth/logout/}
      $\rightarrow$ Destroy session $\rightarrow$ Clear cookie\\[0.3cm]
    \textit{[ Insert sequence diagram image here ]}
  \vspace{1.5cm}}}
  \caption{SD01 --- Django Session Authentication Flow}
  \label{fig:sd01_auth}
\end{figure}

\subsubsection{CSRF Protection}

Cross-Site Request Forgery (CSRF) protection is enforced on all state-changing
requests (\texttt{POST}, \texttt{PUT}, \texttt{DELETE}). Django generates a
\texttt{csrftoken} cookie that is readable by JavaScript. The Axios HTTP client
in the React frontend reads this cookie and attaches its value in the
\texttt{X-CSRFToken} header of every mutating request. Django's
\texttt{CsrfViewMiddleware} then verifies this token server-side before
processing the request body.

The Axios configuration includes two key settings:

\begin{itemize}
  \item \texttt{withCredentials:~true} --- instructs the browser to include the
    \texttt{sessionid} cookie on every cross-origin API request.
  \item A request interceptor that reads the \texttt{csrftoken} cookie and
    attaches it to the \texttt{X-CSRFToken} header on \texttt{POST},
    \texttt{PUT}, and \texttt{DELETE} requests.
\end{itemize}

\subsubsection{Custom User Model}

Django's default \texttt{User} model was extended with a \texttt{CustomUser}
model, adding a \texttt{role} field (\texttt{admin} or \texttt{user}) and a
\texttt{temporary\_password} boolean flag.

\begin{table}[H]
\centering
\caption{CustomUser Model Fields}
\label{tab:custom_user}
\renewcommand{\arraystretch}{1.5}
\begin{tabular}{|l|l|p{7cm}|}
\hline
\rowcolor{blue!20}
\textbf{Field} & \textbf{Type} & \textbf{Description} \\
\hline
\texttt{username} & \texttt{CharField} &
  Unique login identifier for each MARETAP platform user. \\
\hline
\rowcolor{gray!10}
\texttt{email} & \texttt{EmailField} &
  Email address used for temporary password delivery. \\
\hline
\texttt{role} & \texttt{CharField} (choices) &
  Either \texttt{admin} or \texttt{user}; governs all access control
  decisions. \\
\hline
\rowcolor{gray!10}
\texttt{temporary\_password} & \texttt{BooleanField} &
  Set to \texttt{True} when an administrator creates an account; cleared
  after the user changes their password on first login. \\
\hline
\texttt{first\_name} & \texttt{CharField} &
  Displayed in the profile section of the sidebar. \\
\hline
\rowcolor{gray!10}
\texttt{last\_name} & \texttt{CharField} &
  Displayed alongside the first name in the user interface. \\
\hline
\texttt{password} & \texttt{CharField} &
  Password hash stored using PBKDF2 + SHA-256 with a per-user salt.
  Inherited from Django's \texttt{AbstractUser}. \\
\hline
\end{tabular}
\end{table}

\subsection{Role-Based Access Control}

\subsubsection{Backend Permission Classes}

Two custom DRF permission classes are defined in the \texttt{accounts}
application. Both verify an active Django session before evaluating the user's
role:

\begin{itemize}
  \item \texttt{IsAdminRole} --- grants access only to authenticated users whose
    \texttt{role} equals \texttt{"admin"}. Applied to user management,
    Power~BI configuration, and file ingestion endpoints.
  \item \texttt{IsAuthenticatedWithAnyRole} --- grants access to any
    authenticated user regardless of role. Applied to KPI, chatbot, and
    report-viewing endpoints.
\end{itemize}

\subsubsection{Frontend Route Guards}

The \texttt{ProtectedRoute} higher-order component calls
\texttt{GET /api/auth/me/} on mount, which reads the active session and returns
the current user's role. If the required role is absent, the user is redirected
to a 403 error page.

\subsection{REST API Endpoints}

All endpoints require an active Django session (\texttt{sessionid} cookie).
Write operations additionally require a valid CSRF token in the
\texttt{X-CSRFToken} header.

% ── SEQUENCE DIAGRAM: KPI Fetching ───────────────────────────
\begin{figure}[H]
  \centering
  \fbox{\parbox{0.9\textwidth}{\centering\vspace{1.5cm}
    \textbf{[ SEQUENCE DIAGRAM ]}\\[0.3cm]
    SD02 --- KPI Dashboard Data Flow\\[0.2cm]
    React (\texttt{sessionid} cookie)
      $\rightarrow$ \texttt{GET /api/dashboard/kpis/}\\
    $\rightarrow$ Django: verify session
      $\rightarrow$ ORM query $\rightarrow$ SQL Server DWH\\
    $\rightarrow$ JSON response $\rightarrow$ Chart.js rendering\\[0.3cm]
    \textit{[ Insert sequence diagram image here ]}
  \vspace{1.5cm}}}
  \caption{SD02 --- KPI Data Fetching Sequence Diagram}
  \label{fig:sd02_kpi}
\end{figure}

\begin{table}[H]
\centering
\caption{EZZAOUIA REST API --- Complete Endpoint Reference}
\label{tab:api_endpoints}
\renewcommand{\arraystretch}{1.5}
\begin{tabular}{|l|l|l|p{4.5cm}|}
\hline
\rowcolor{blue!20}
\textbf{Method} & \textbf{Endpoint} & \textbf{Access} & \textbf{Description} \\
\hline
POST & \texttt{/api/auth/login/} & Public &
  Validates credentials, creates a Django session, and sets the
  \texttt{sessionid} cookie. \\
\hline
\rowcolor{gray!10}
POST & \texttt{/api/auth/logout/} & Authenticated &
  Destroys the server-side session and clears the \texttt{sessionid}
  cookie. \\
\hline
GET & \texttt{/api/auth/me/} & Authenticated &
  Returns the current user's information and role from the active
  session. \\
\hline
\rowcolor{gray!10}
GET & \texttt{/api/dashboard/kpis/} & Authenticated &
  Returns BOPD, total gas, total water, and field-level summary KPIs. \\
\hline
GET & \texttt{/api/dashboard/bsw/} & Authenticated &
  Returns average BSW per well for the current period. \\
\hline
\rowcolor{gray!10}
GET & \texttt{/api/dashboard/gor/} & Authenticated &
  Returns average GOR per well for the current period. \\
\hline
GET & \texttt{/api/dashboard/top-wells/} & Authenticated &
  Returns the five wells with the highest daily oil production. \\
\hline
\rowcolor{gray!10}
GET & \texttt{/api/dashboard/monthly-trend/} & Authenticated &
  Returns monthly aggregated oil production for the current year. \\
\hline
GET & \texttt{/api/dashboard/tank-levels/} & Authenticated &
  Returns current stock volumes and safety margins per tank. \\
\hline
\rowcolor{gray!10}
GET & \texttt{/api/reports/} & Authenticated &
  Returns the list of Power~BI reports with their embed URLs. \\
\hline
GET/POST & \texttt{/api/users/} & Admin only &
  Lists all users or creates a new user account. \\
\hline
\rowcolor{gray!10}
PUT/DELETE & \texttt{/api/users/\{id\}/} & Admin only &
  Updates or deletes a specific user account. \\
\hline
POST & \texttt{/api/chatbot/query/} & Authenticated &
  Submits a natural language query to the RAG pipeline. \\
\hline
\rowcolor{gray!10}
POST & \texttt{/api/chatbot/ingest/} & Admin only &
  Triggers document ingestion and ChromaDB indexing via a Celery task. \\
\hline
GET & \texttt{/api/chatbot/ingest/status/\{task\_id\}/} & Authenticated &
  Returns the current status of a Celery ingestion task
  (\textit{Pending}, \textit{Processing}, \textit{Success}, or
  \textit{Failed}). \\
\hline
\end{tabular}
\end{table}

\subsection{KPI Computation Logic}

\subsubsection{BOPD}
Computed as the sum of \texttt{DailyOilPerWellSTBD} across all active wells for
the most recent date, using Django ORM's \texttt{aggregate(Sum(...))}.

\subsubsection{BSW per Well}
Per-well average over the last 30~days, obtained by chaining
\texttt{values("well\_\_WellCode")} with
\texttt{annotate(avg\_bsw=Avg("dimwellstatus\_\_BSW"))}, ordered
descending.

\subsubsection{Monthly Production Trend}
\texttt{TruncMonth} groups records by calendar month;
\texttt{annotate(monthly\_oil=Sum("DailyOilPerWellSTBD"))} produces the
12-month ordered series used by the Chart.js line chart.

\subsection{User Management API}

\subsubsection{Account Creation with Temporary Password}

When an administrator creates a new user, the backend generates a random
12-character password using Python's \texttt{secrets} module, sets
\texttt{temporary\_password = True}, and dispatches a welcome email. The API
response includes an \texttt{email\_sent} boolean indicating whether SMTP
delivery succeeded.

% ── SEQUENCE DIAGRAM: User Creation ──────────────────────────
\begin{figure}[H]
  \centering
  \fbox{\parbox{0.9\textwidth}{\centering\vspace{1.5cm}
    \textbf{[ SEQUENCE DIAGRAM ]}\\[0.3cm]
    SD03 --- User Creation Flow\\[0.2cm]
    Admin
      $\rightarrow$ \texttt{POST /api/users/}
        (\texttt{sessionid} + \texttt{X-CSRFToken})\\
    $\rightarrow$ Django: verify session \& CSRF
      $\rightarrow$ Generate password (\texttt{secrets})\\
    $\rightarrow$ SMTP email
      $\rightarrow$ \texttt{\{email\_sent: true/false\}}\\
    $\rightarrow$ Persistent password modal (until manually closed)\\[0.3cm]
    \textit{[ Insert sequence diagram image here ]}
  \vspace{1.5cm}}}
  \caption{SD03 --- User Management Sequence Diagram}
  \label{fig:sd03_user}
\end{figure}

\subsubsection{SMTP Configuration}

Email is sent via Django's \texttt{EmailMessage} over SMTP port~587
(STARTTLS) to Microsoft~365. By default, Microsoft~365 disables SMTP
Authentication at the tenant level; enabling it requires an Office~365
administrator to activate \textit{Authenticated SMTP} in the Microsoft~365
Admin Center or via PowerShell. This is a tenant-level configuration action
and cannot be addressed within the Django application code.

\end{sloppypar}

% ─────────────────────────────────────────────────────────────
\section{Frontend Implementation}
\begin{sloppypar}

\subsection{React Application Structure}

\begin{table}[H]
\centering
\caption{React Frontend Directory Structure}
\label{tab:react_structure}
\renewcommand{\arraystretch}{1.5}
\begin{tabular}{|l|p{9cm}|}
\hline
\rowcolor{blue!20}
\textbf{Directory / File} & \textbf{Contents} \\
\hline
\texttt{src/components/} &
  \texttt{Sidebar.jsx}, \texttt{ProtectedRoute.jsx}, \texttt{KPICard.jsx},
  \texttt{PasswordModal.jsx}. \\
\hline
\rowcolor{gray!10}
\texttt{src/pages/} &
  \texttt{LoginPage.jsx}, \texttt{DashboardPage.jsx},
  \texttt{ReportsPage.jsx}, \texttt{ChatbotPage.jsx},
  \texttt{UsersPage.jsx}, \texttt{FilesPage.jsx}. \\
\hline
\texttt{src/api/} &
  Axios instance configured with \texttt{withCredentials:~true} and a
  CSRF request interceptor. \\
\hline
\rowcolor{gray!10}
\texttt{src/context/} &
  \texttt{AuthContext.jsx} --- fetches user information and role from
  \texttt{/api/auth/me/} on application mount. \\
\hline
\texttt{src/hooks/} &
  \texttt{useKPIs.js}, \texttt{useReports.js}, \texttt{useUsers.js}. \\
\hline
\rowcolor{gray!10}
\texttt{src/App.jsx} &
  Root routing tree with \texttt{ProtectedRoute} guards applied per
  role. \\
\hline
\end{tabular}
\end{table}

\subsection{Authentication Context and Axios Configuration}

On application mount, \texttt{AuthContext} calls
\texttt{GET /api/auth/me/}. If the browser holds a valid \texttt{sessionid}
cookie, Django returns the user's information and role, restoring the
authenticated session seamlessly without requiring the user to log in
again.

Axios is configured with two key settings:

\begin{itemize}
  \item \texttt{withCredentials:~true} --- the browser includes the
    \texttt{sessionid} cookie on every cross-origin request automatically.
  \item A \textbf{request interceptor} that reads the \texttt{csrftoken} cookie
    and attaches its value to the \texttt{X-CSRFToken} header on
    \texttt{POST}, \texttt{PUT}, and \texttt{DELETE} requests.
\end{itemize}

\subsection{Role-Gated Sidebar Navigation}

\begin{table}[H]
\centering
\caption{Sidebar Navigation Items by Role}
\label{tab:sidebar_roles}
\renewcommand{\arraystretch}{1.5}
\begin{tabular}{|l|c|c|}
\hline
\rowcolor{blue!20}
\textbf{Navigation Item} & \textbf{Admin} & \textbf{User} \\
\hline
Dashboard / Overview      & \checkmark & \checkmark \\
\hline
\rowcolor{gray!10}
Power~BI Reports          & \checkmark & \checkmark \\
\hline
AI Chatbot                & \checkmark & \checkmark \\
\hline
\rowcolor{gray!10}
User Management           & \checkmark & $\times$ \\
\hline
File Ingestion            & \checkmark & $\times$ \\
\hline
\rowcolor{gray!10}
Administration            & \checkmark & $\times$ \\
\hline
\end{tabular}
\end{table}

The Logout button calls \texttt{POST /api/auth/logout/} (with the CSRF token
in the header), clears the \texttt{AuthContext} state, and redirects the user
to the login page.

\begin{figure}[H]
    \centering
    \includegraphics[width=0.25\textwidth]{images/sidebar_admin.png}
    \caption{EZZAOUIA Sidebar --- Admin View with Full Navigation}
    \label{fig:sidebar_admin}
\end{figure}

\begin{figure}[H]
    \centering
    \includegraphics[width=0.25\textwidth]{images/sidebar_user.png}
    \caption{EZZAOUIA Sidebar --- Standard User View with Restricted Navigation}
    \label{fig:sidebar_user}
\end{figure}

\subsection{KPI Dashboard Page}

The KPI Dashboard is the default landing page after authentication. The first
row presents four KPI cards (BOPD, total gas, BSW, GOR) with colour-coded
thresholds. The second row displays a horizontal bar chart of the top~5 wells
by production, a line chart of the 12-month production trend, and a doughnut
chart showing the production breakdown by artificial lift method.

\begin{figure}[H]
    \centering
    \includegraphics[width=0.95\textwidth]{images/kpi_dashboard.png}
    \caption{EZZAOUIA KPI Dashboard --- Production Overview Page}
    \label{fig:kpi_dashboard}
\end{figure}

\subsection{Power~BI Embedded Reports Page}

Four production dashboards are embedded via \texttt{<iframe>} using the
following URL pattern:

\begin{verbatim}
https://app.powerbi.com/reportEmbed?reportId={id}
&autoAuth=true&ctid={tenant_id}
\end{verbatim}

\begin{figure}[H]
    \centering
    \includegraphics[width=0.95\textwidth]{images/powerbi_embedded.png}
    \caption{Power~BI Overview Dashboard Embedded in the EZZAOUIA Platform}
    \label{fig:powerbi_embedded}
\end{figure}

\subsection{User Management Interface}

The user management page provides the following behaviour following account
creation:

\begin{itemize}
  \item If \texttt{email\_sent === true}: a persistent password modal with a
    copy button is displayed until the administrator manually dismisses it.
  \item If \texttt{email\_sent === false}: the same modal is shown together
    with a persistent amber warning banner, prompting the administrator to
    communicate the credentials manually.
\end{itemize}

\begin{figure}[H]
    \centering
    \includegraphics[width=0.85\textwidth]{images/user_management.png}
    \caption{User Management Page --- Account List and Creation Interface}
    \label{fig:user_management}
\end{figure}

\begin{figure}[H]
    \centering
    \includegraphics[width=0.65\textwidth]{images/password_modal.png}
    \caption{Temporary Password Modal with Email Failure Warning Banner}
    \label{fig:password_modal}
\end{figure}

\end{sloppypar}

% ─────────────────────────────────────────────────────────────
\section{AI Chatbot Module (RAG Architecture)}
\begin{sloppypar}

\subsection{Motivation and Design Choice}

The AI chatbot allows MARETAP's engineers to ask natural language questions about
production data, well performance, and uploaded technical documents. The platform
implements a \textbf{Retrieval-Augmented Generation (RAG)} architecture that
grounds every response in content retrieved from the company's own document
corpus, without requiring expensive model fine-tuning or cloud API access.

\noindent\textit{Note: the RAG chatbot is presented here as an integrated module
of the platform. A dedicated chapter provides a detailed treatment of the RAG
pipeline design, evaluation, and implementation.}

\subsection{RAG Pipeline Architecture}

% ── SEQUENCE DIAGRAM: RAG ────────────────────────────────────
\begin{figure}[H]
  \centering
  \fbox{\parbox{0.9\textwidth}{\centering\vspace{1.5cm}
    \textbf{[ SEQUENCE DIAGRAM ]}\\[0.3cm]
    SD04 --- RAG Pipeline --- Two Phases\\[0.2cm]
    \textbf{Ingestion:} Upload
      $\rightarrow$ Celery task $\rightarrow$ Extract text\\
    $\rightarrow$ Chunk (500~tokens / 50-token overlap)
      $\rightarrow$ Embed $\rightarrow$ ChromaDB\\[0.2cm]
    \textbf{Query:} Question
      $\rightarrow$ Embed $\rightarrow$ ChromaDB ($k=5$)\\
    $\rightarrow$ Context assembly
      $\rightarrow$ Ollama generate $\rightarrow$ Response\\[0.3cm]
    \textit{[ Insert sequence diagram image here ]}
  \vspace{1.5cm}}}
  \caption{SD04 --- RAG Pipeline Architecture Sequence Diagram}
  \label{fig:sd04_rag}
\end{figure}

\subsubsection{Ingestion Phase}

\begin{itemize}
  \item \textbf{Text extraction:} \texttt{PyMuPDF} for PDF files;
    \texttt{python-docx} for DOCX files.
  \item \textbf{Chunking:} 500-token windows with a 50-token overlap (sliding
    window strategy).
  \item \textbf{Embedding:} Ollama's embedding endpoint converts each chunk
    into a dense vector representation.
  \item \textbf{Indexing:} Vectors are stored in ChromaDB with source metadata
    for attribution at query time.
\end{itemize}

\subsubsection{Query Phase}

\begin{itemize}
  \item \textbf{Query embedding:} The user's question is embedded using the
    same model as ingestion.
  \item \textbf{Semantic retrieval:} ChromaDB performs a nearest-neighbour
    search and returns the top $k=5$ chunks.
  \item \textbf{Context assembly:} Retrieved chunks are concatenated with
    source attribution and inserted into the prompt.
  \item \textbf{Response generation:} The augmented prompt is submitted to
    Ollama; the grounded response is returned to the frontend.
\end{itemize}

\begin{figure}[H]
    \centering
    \includegraphics[width=0.95\textwidth]{images/rag_pipeline.png}
    \caption{RAG Pipeline Architecture --- Ingestion and Query Phases}
    \label{fig:rag_pipeline}
\end{figure}

\begin{figure}[H]
    \centering
    \includegraphics[width=0.85\textwidth]{images/chatbot_interface.png}
    \caption{EZZAOUIA AI Chatbot Interface}
    \label{fig:chatbot}
\end{figure}

\end{sloppypar}

% ─────────────────────────────────────────────────────────────
\section{File Ingestion Module}
\begin{sloppypar}

\subsection{Asynchronous Processing with Celery}

File ingestion is implemented as an asynchronous Celery task. The upload
endpoint returns a task~ID immediately upon receiving the request; the Celery
worker then processes the document in the background, dispatching its work via
the Redis message broker. This design ensures that large files and long
extraction operations never block the main Django request thread.

% ── SEQUENCE DIAGRAM: Ingestion ──────────────────────────────
\begin{figure}[H]
  \centering
  \fbox{\parbox{0.9\textwidth}{\centering\vspace{1.5cm}
    \textbf{[ SEQUENCE DIAGRAM ]}\\[0.3cm]
    SD05 --- Asynchronous File Ingestion Flow\\[0.2cm]
    Admin
      $\rightarrow$ \texttt{POST /api/chatbot/ingest/}
        (\texttt{sessionid} + \texttt{X-CSRFToken})\\
    $\rightarrow$ Task~ID returned immediately
      $\rightarrow$ Redis queue\\
    $\rightarrow$ Celery worker
      $\rightarrow$ ChromaDB indexing\\
    $\rightarrow$ Frontend polls
      \texttt{/status/\{task\_id\}/} every 2~s\\[0.3cm]
    \textit{[ Insert sequence diagram image here ]}
  \vspace{1.5cm}}}
  \caption{SD05 --- Asynchronous File Ingestion Sequence Diagram}
  \label{fig:sd05_ingestion}
\end{figure}

\subsection{Ingestion Status Tracking}

The frontend polls
\texttt{/api/chatbot/ingest/status/\{task\_id\}/} every two seconds and
displays one of four states: \textit{Pending}, \textit{Processing},
\textit{Success}, or \textit{Failed}.

\begin{figure}[H]
    \centering
    \includegraphics[width=0.85\textwidth]{images/file_ingestion.png}
    \caption{File Ingestion Page --- Document Upload and Indexing Interface}
    \label{fig:file_ingestion}
\end{figure}

\end{sloppypar}

% ─────────────────────────────────────────────────────────────
\section{Security Architecture}
\begin{sloppypar}

\subsection{Authentication and Session Security}

The platform's security rests on three complementary layers:

\begin{enumerate}
  \item \textbf{Django Session Authentication:} Every endpoint verifies a valid
    \texttt{sessionid} cookie. Sessions are stored server-side --- invalidating
    a session immediately terminates access, even if the browser still holds
    the cookie. The \texttt{sessionid} cookie is \texttt{HttpOnly}, preventing
    XSS-based session hijacking.

  \item \textbf{CSRF Token Verification:} Django's \texttt{CsrfViewMiddleware}
    is active on all \texttt{POST}, \texttt{PUT}, and \texttt{DELETE}
    endpoints. Requests that are missing a valid \texttt{X-CSRFToken} header
    are rejected with HTTP~403 before any business logic is executed. This
    prevents cross-site request forgery attacks, where a malicious third-party
    page might attempt to trigger actions using the victim's session cookie.

  \item \textbf{Role-Based Access Control:} Enforced at the API layer through
    DRF permission classes and at the frontend through \texttt{ProtectedRoute}
    guards, ensuring that unauthorised requests are rejected before any
    sensitive data is exposed.
\end{enumerate}

\subsection{Cross-Origin Resource Sharing (CORS)}

\texttt{django-cors-headers} whitelists the React frontend's origin.
\texttt{CORS\_ALLOW\_CREDENTIALS = True} permits the browser to include the
\texttt{sessionid} cookie on cross-origin requests. In production, only
IP addresses belonging to MARETAP's internal network are whitelisted.

\subsection{Password Security}

Passwords are stored using PBKDF2 + SHA-256 with a per-user salt, the default
Django password hashing backend. Temporary passwords are generated via
\texttt{secrets.token\_urlsafe()}. The \texttt{temporary\_password} flag
enforces a mandatory password change on first login.

\subsection{SQL Server Access Security}

The \texttt{stage} service account holds read-only permissions on the Data
Warehouse. Connection encryption is enforced via
\texttt{encrypt=true;trustServerCertificate=true} in the ODBC connection
string.

\end{sloppypar}

% ─────────────────────────────────────────────────────────────
\section{Containerised Deployment}
\begin{sloppypar}

\subsection{Docker Compose Architecture}

The EZZAOUIA platform is packaged as a four-service Docker Compose application.
Two additional dependencies --- ChromaDB and Ollama --- run on the host machine
outside the Docker network and are reached by the containers via the host
network interface.

\begin{table}[H]
\centering
\caption{Docker Compose Services of the EZZAOUIA Platform}
\label{tab:docker_services}
\renewcommand{\arraystretch}{1.5}
\begin{tabular}{|l|l|p{7cm}|}
\hline
\rowcolor{blue!20}
\textbf{Service} & \textbf{Image} & \textbf{Role} \\
\hline
\texttt{web} & \texttt{ezzaouia\_platform-web} &
  Django REST API on port~8000. Handles session authentication, CSRF
  protection, KPI computation, and all platform business logic. \\
\hline
\rowcolor{gray!10}
\texttt{celery} & \texttt{ezzaouia\_platform-celery} &
  Asynchronous Celery worker. Consumes document ingestion tasks
  dispatched by Django to the Redis broker queue. No exposed port. \\
\hline
\texttt{celery-beat} & \texttt{ezzaouia\_platform-celery-beat} &
  Celery Beat periodic task scheduler. Triggers time-based background
  jobs. No exposed port. \\
\hline
\rowcolor{gray!10}
\texttt{redis} & \texttt{redis:7-alpine} &
  Celery message broker on port~6379. Runs inside the Docker Compose
  network; no external dependency required. \\
\hline
\end{tabular}
\end{table}

\noindent\textit{Note: ChromaDB (vector store) and Ollama (LLM inference
server) are not containerised. They run as host-machine processes and are
accessed by the \texttt{web} and \texttt{celery} containers via the host
network. ChromaDB persists its data to \texttt{./chroma\_db/} on the host
filesystem.}

\subsection{Build and Deployment Procedure}

Every code change requires a full image rebuild before the updated containers
are started:

\begin{verbatim}
docker compose down && docker compose up --build
\end{verbatim}

\subsection{Environment Configuration}

Sensitive values --- database credentials, Django \texttt{SECRET\_KEY}, SMTP
password, Power~BI tenant ID, and session parameters --- are injected at
runtime via a \texttt{.env} file that is excluded from version control.

\end{sloppypar}

% ─────────────────────────────────────────────────────────────
\section{Platform Interface Overview}
\begin{sloppypar}

\subsection{Login Page}

\begin{figure}[H]
    \centering
    \includegraphics[width=0.85\textwidth]{images/login_page.png}
    \caption{EZZAOUIA Login Page}
    \label{fig:login}
\end{figure}

\subsection{KPI Dashboard}

The default authenticated landing page, aggregating BOPD, total gas production,
BSW, GOR, the top-performing wells, the 12-month production trend, and the
production breakdown by artificial lift method into a single real-time
view.

\subsection{Power~BI Reports Page}

Tab-based access to all four embedded Power~BI dashboards (Overview, Sales,
Tanks, and Well Health) with full-screen rendering and no external navigation
required.

\subsection{Profile Management}

All authenticated users can view and update their personal information
(first name, last name, and email) and change their password via a dedicated
profile page. Users who are still carrying a temporary password are prompted
to change it on first access.

\subsection{Administration Interface}

Visible only to administrators: user management, file ingestion, and system
configuration are consolidated in one place, eliminating the need for direct
database access.

\end{sloppypar}

% ─────────────────────────────────────────────────────────────
\section{Conclusion}
\begin{sloppypar}

This chapter presented the complete design, implementation, and deployment of the
EZZAOUIA web platform --- the central contribution of this internship project.
Over the course of two development sprints, a full-stack web application was
constructed that transforms the structured data of the MARETAP Data Warehouse
into an accessible, secure, and intelligent operational tool for
MARETAP~S.A.'s production teams.

The Django backend provides a clean REST API protected by Django's native
session-based authentication and CSRF token verification --- a deliberate
architectural choice well-suited to the on-premise deployment context of
MARETAP's internal network, where server-side session management is both simpler
and equally secure compared to token-based alternatives. The backend exposes KPI
data, manages user accounts, orchestrates document ingestion through Celery, and
serves Power~BI report metadata to the frontend.

The React frontend delivers a role-gated, single-page user experience built
around an adaptive sidebar, a real-time KPI dashboard powered by Chart.js, an
embedded Power~BI reporting module, an AI chatbot grounded in a RAG architecture
using Ollama and ChromaDB, and a complete user management interface with
temporary password handling and SMTP email integration.

The entire platform is packaged as a \textbf{four-service} Docker Compose
application --- \texttt{web} (Django), \texttt{celery} (asynchronous worker),
\texttt{celery-beat} (periodic scheduler), and \texttt{redis} (message broker)
--- with ChromaDB and Ollama running on the host machine outside Docker. This
architecture ensures reproducible deployment across MARETAP's on-premise
environment with no cloud dependency.

The security architecture --- built on Django session authentication, CSRF token
enforcement, role-based access control at both the API and frontend levels, and
read-only Data Warehouse access through a dedicated service account --- ensures
that the platform meets the confidentiality and integrity requirements of an oil
production environment. Together, the Data Warehouse constructed in the previous
chapter and the EZZAOUIA platform presented here constitute a complete,
production-ready business intelligence and AI-assisted monitoring solution for
the Ezzaouia field.

\end{sloppypar}
